\documentclass[11pt,a4paper]{article}
\usepackage[margin=23mm]{geometry}
\usepackage[T1]{fontenc}
\usepackage{lmodern}
\usepackage{microtype}
\usepackage{xcolor,tabularx,booktabs,array,enumitem,hyperref,fancyhdr}
\definecolor{navy}{HTML}{163654}
\definecolor{teal}{HTML}{147C84}
\hypersetup{colorlinks=true,urlcolor=teal,linkcolor=navy}
\setlength{\parindent}{0pt}\setlength{\parskip}{5pt}
\pagestyle{fancy}\fancyhf{}\lhead{\small Wearable Athlete Monitoring System}\rhead{\small MNE1 -- Week 1}\cfoot{\thepage}
\newcommand{\sect}[1]{\vspace{6pt}{\color{navy}\large\bfseries #1}\par\vspace{2pt}}
\newcommand{\source}[1]{{\footnotesize\color{teal}Source: #1}}
\begin{document}
{\color{navy}\LARGE\bfseries Sensor Interfaces and Electrical Compatibility}\par
{\large Wearable Athlete Monitoring System -- Week 1}\par\vspace{9pt}
\begin{tabular}{@{}ll@{}}\textbf{Author:} & Nour El Islam Zougham \\
\textbf{Responsibility:} & MNE1 -- Electronics \& Hardware\\
\textbf{Date:} & 20 November 2026\\
\textbf{Scope:} & Literature and preliminary compatibility study; no bench validation\\
\end{tabular}
\vspace{4pt}\hrule\vspace{8pt}
\textbf{Purpose.} This report develops the author's handwritten research on two available sensors: an MPU6050 motion sensor and a DS18B20-based temperature module (visually consistent with KY-001). Technical corrections are incorporated, and assumptions requiring physical verification are clearly identified. No measurement or successful connection is claimed.

\sect{1. MPU6050 motion-sensing module}
\textbf{Function and interface.} The MPU6050 is a six-axis inertial measurement IC combining a three-axis accelerometer and a three-axis gyroscope. It uses \textbf{I\textsuperscript{2}C} for communication. \textbf{SDA} is the bidirectional serial data line; \textbf{SCL} carries the clock generated by the bus controller. The Arduino typically acts as the controller and the sensor as a peripheral. Pull-up resistors hold the bus lines high when no device is pulling them low; they must pull up to a \emph{safe and compatible logic rail}, not automatically to 5 V.

\textbf{Device address.} The MPU6050 normally has the 7-bit I\textsuperscript{2}C address \texttt{0x68} when AD0 is low, and \texttt{0x69} when AD0 is high. An I\textsuperscript{2}C scanner may help confirm a responsive device during a \emph{later} powered test; scanning is not a prerequisite to understanding its specifications.

\textbf{Electrical specifications and corrections.}\par
\begin{tabularx}{\linewidth}{@{}p{39mm}X@{}}\toprule
IC VDD supply & \textbf{2.375--3.46 V}; 3.3 V is a common nominal supply. This describes the \emph{IC}, not necessarily the breakout connector's permissible VCC.\\
IC logic domain & The MPU6050 has a VLOGIC domain; SDA/SCL tolerances depend on the IC's logic and supply rails. Its pins must \emph{not} be assumed 5-V tolerant.\\
Current consumption & \textbf{3.6 mA} is approximately a \emph{typical gyroscope-only} operating current, not the maximum current of the device. Typical current with gyroscope, accelerometer and DMP enabled is about \textbf{3.9 mA}; current varies by mode.\\
Module supply and level shifting & \textbf{Not verified.} A breakout may include a regulator and/or pull-ups; regulator presence alone does not establish safe 5-V logic compatibility.\\
\bottomrule\end{tabularx}

\textbf{Arduino Uno assessment.} A conventional Arduino Uno uses 5-V logic. It can communicate over I\textsuperscript{2}C with a correctly powered MPU6050 breakout only if the SDA/SCL voltage levels are compatible. The handwritten estimate of approximately \textbf{3.15 V} from interacting pull-up resistors is \emph{not a verified measurement or a general design guarantee}. Enabling 5-V internal pull-ups can expose the sensor to an unsafe voltage; confirm board resistor connections and, where necessary, use suitable bidirectional I\textsuperscript{2}C level shifting. \textbf{Never connect AD0 directly to 5 V} unless the particular board explicitly makes that safe.

\source{MPU-6000/MPU-6050 Product Specification, Rev. 3.4: \url{https://hw101.tbs1.de/mpu6050/doc/mpu6050.pdf}}

\sect{2. DS18B20 temperature-sensing module}
\textbf{Identification and interface.} The pictured module resembles a \textbf{KY-001} board with a TO-92-style \textbf{DS18B20} digital thermometer. Confirm the actual component marking before treating the identification as certain. Unlike the MPU6050, the DS18B20 uses a \textbf{1-Wire} bus: one bidirectional digital data line, \textbf{DQ}, plus power and ground when operated with external supply.

\textbf{Pin distinction.} With the bare TO-92 package's flat face toward the viewer and leads pointing downward, its pins are \textbf{left: GND; middle: DQ; right: VDD}. The \emph{module's} three-pin header, as shown in the earlier board photograph, is labeled \textbf{-- (GND), + (VCC), OUT (data)}. Bare-IC lead order and PCB-header order must not be confused.

\textbf{Pull-up resistor.} A 1-Wire bus needs a pull-up on DQ. A common reference value for a basic, short, externally powered bus is \textbf{4.7 k\(\Omega\)}. However, this is \emph{not always an extra breadboard part}: some KY-001-style modules have a pull-up already fitted. Inspect the particular PCB and resistance markings (a marking \texttt{103}, for example, denotes 10 k\(\Omega\)) before adding another resistor. The sentence ````it will not work without adding a 4.7-k\(\Omega\) resistor'' is therefore too absolute.

\textbf{Electrical supply and Uno compatibility.} The DS18B20 IC operates from \textbf{3.0--5.5 V}, rather than only 3.3--5 V. It can normally operate on a 5-V Arduino Uno system with an appropriately wired data bus and pull-up. The data connection is to a \textbf{digital GPIO pin}, not the Uno's I\textsuperscript{2}C SDA/SCL pins. A specific Uno digital pin (for example D2) may be selected in software; no special pin named ````D0'' or analog input is required. The handwritten statement of ````100\% compatible'' should be replaced with \emph{electrically compatible in a correctly configured 5-V circuit}, subject to verifying the exact module.

\source{Analog Devices, DS18B20 Datasheet: \url{https://www.analog.com/media/en/technical-documentation/data-sheets/DS18B20.pdf}}

\sect{3. Consolidated Week 1 interface shortlist}
\small
\begin{tabularx}{\linewidth}{@{}p{29mm}p{29mm}p{24mm}X@{}}\toprule
\textbf{Sensor / status} & \textbf{IC operating supply} & \textbf{Interface} & \textbf{Open electrical checks}\\\midrule
MPU6050 (owned) & 2.375--3.46 V & I\textsuperscript{2}C; SDA/SCL & Exact breakout model; connector VCC rating; pull-up rail; 5-V Uno logic interface.\\
DS18B20 on KY-001-style board (owned; marking unconfirmed) & 3.0--5.5 V & 1-Wire; DQ/OUT & Confirm IC marking, fitted pull-up value, header labels and module revision.\\
\bottomrule
\end{tabularx}\normalsize

\textbf{Scope boundary.} This report documents the two owned sensors only. It does not claim the Week 1 multimeter exercise, voltage-domain sketch, sourcing inquiries for additional candidates, or team decisions have been completed. Nor does it constitute a wiring schematic, successful bus test, or demonstration of accurate skin/core-temperature measurement.

\vspace{7pt}\hrule\vspace{3pt}
{\footnotesize\textbf{Authorship note.} Original research and handwritten notes: Nour El Islam Zougham. Technical proofreading, editing, typesetting, and PDF compilation assisted by ChatGPT.}
\end{document}